\section{Computing the full height in polynomial space}\label{fg:sec:space}
Fixed thresholds use a constant number of Stienstra coefficients. Infinite
height requires a finite stopping argument. A finite operator from the
quasi-$F$-splitting criterion supplies one, and its entries can be computed
from a circuit without storing its exponentially large matrix.

We use the following theorem in its established geometric form.
\citet[Theorem~C, p.~4, and Theorem~5.8]{fg-kty-one} give a Fedder type
criterion for quasi-$F$-split height. For smooth Calabi--Yau varieties this
height agrees with Artin--Mazur height by
\citet[Theorem~4.5, pp.~9--10 of the preprint]{fg-yobuko}.
That source assumes an algebraically closed ground field. Apply its
comparison after extension to $\overline{\Fp}$: the first nonzero
coefficient of a formal group's $p$-series is unchanged by field
extension, and the coefficient recurrence below for an equation over
$\Fp$ stays in $\Fp$. Its first nonzero index is therefore unchanged
as well; see also \citet[Theorem~5.13]{fg-kty-one}.
The concrete monomial formula and finite operator
are given by \citet[Proposition~3.2, Remark~3.4 and Algorithms~2--4,
pp.~5--7]{fg-bgp}. We recall the exact criterion and prove its complexity
consequences for circuit input.

Let $F\in\Fp[x_1,\ldots,x_d]$ be homogeneous of degree $d\ge3$, and
assume its projective hypersurface is smooth and geometrically connected.
Put
\[
 D=d(p-1),\qquad g=F^{p-1},\qquad
 V=\Fp[x_1,\ldots,x_d]_D,\qquad
 B=\dim V=\binom{pd-1}{d-1}.
\]
For the collected monomial expansion $g=\sum_\beta c_\beta x^\beta$,
choose integer representatives $\widehat c_\beta\in\{0,\ldots,p-1\}$,
set $G=\sum_\beta\widehat c_\beta x^\beta$, and define
\begin{equation}\label{fg:eq:delta}
 \Delta=\left(\frac{G^p-\sum_\beta
                    (\widehat c_\beta x^\beta)^p}{p}\right)\bmod p.
\end{equation}
The numerator is coefficientwise divisible by $p$. Changing the integer
representatives changes both $p$-th powers by multiples of $p^2$, so
the result is well defined. Define the $\Fp$-linear extractor $u$ by
\[
 u(x^\gamma)=
 \begin{cases}
 x^{(\gamma-(p-1)\mathbf1)/p},
       &\gamma\equiv(p-1)\mathbf1\pmod p,\\
 0,&\text{otherwise}.
 \end{cases}
\]
Nonnegative exponents in the first case are automatic. Let
$T(a)=u(\Delta a)$ and let $\ell$ select the coefficient with exponent
$\tau=(p-1)\mathbf1$. Since $\deg\Delta=pD$, the operator $T$ preserves
$V$: extracting from degree $(p+1)D$ subtracts $d(p-1)$ and divides by
$p$, returning degree $D$.
The cited criterion is
\begin{equation}\label{fg:eq:finite-operator}
 h(X)=\min\{j+1:\ell(T^jg)\ne0\},
 \qquad \min\varnothing=\infty.
\end{equation}
Indeed, the unique degree-$D$ monomial outside
$(x_1^p,\ldots,x_d^p)$ is $x^\tau$.

The representative convention in \eqref{fg:eq:delta} is compatible with
the Frobenius-image convention for $\Delta_1$ in the criterion.
Grouping monomials with the same exponent residue modulo $p$ changes
the divided difference only by a $p$-th-power polynomial. Replacing
$\Delta$ by $\Delta+H^p$ changes $T(a)$ by $H u(a)$.
For $a\in V$, $u(a)=\ell(a)$ is constant. Thus the two recurrences
agree until their first nonzero target, and agree forever when every
target is zero. This proves the required compatibility for both finite
and infinite height.

\begin{theorem}[A space bound for the full height]\label{fg:thm:space}
Fix $p$. Let $L$ be the complete bit length of a division-free arithmetic
circuit over $\Fp$, with addition, subtraction and multiplication gates,
and with its declared variable count $d$ and all indices encoded in
binary. Promise that it computes a homogeneous degree-$d$ equation in
$d\ge3$ variables defining a smooth geometrically connected projective
hypersurface. There is a deterministic algorithm which outputs its
finite Artin--Mazur height in binary, or outputs infinity, using
\[
 O_p\bigl(L^2\log(L+2)\bigr)
\]
working bits. Every finite height is at most $B=\binom{pd-1}{d-1}$.
The algorithm halts on all syntactically valid inputs and has arbitrary
output outside the promise. Infinite-height decision, and finite-height
or threshold queries with a binary-encoded queried height, therefore
belong to polynomial space.
\end{theorem}

\begin{proof}
Cayley--Hamilton gives a recurrence of order $B$ for
$a_j=\ell(T^jg)$. If $a_0,\ldots,a_{B-1}$ all vanish, every subsequent
term vanishes. Otherwise the first nonzero term has index below $B$.
Together with \eqref{fg:eq:finite-operator}, this proves the height
bound and the finite stopping rule. No characteristic polynomial needs
to be computed.

We give the storage analysis in detail. Every declared variable occurs
on a promised input, by the singular-coordinate-vertex argument in
Proposition~\ref{fg:prop:height-one}; hence $d=O(L)$. A total algorithm
first rejects $d<3$ or a declared dimension larger than its number of
variable occurrences. The bound $B\le2^{pd-1}$ gives $\log B=O_p(d)$. The integer
$B$ can be computed by successive exact divisions in the binomial
formula, using $O_p(d)$ bits. A degree-$D$ basis monomial is represented
by $d$ exponents of size at most $D$, using $O_p(d\log(d+2))$ bits.
Enumerate these weak compositions in a fixed order, with one such
counter, whenever a sum over the basis is required.

First compute one circuit coefficient. For a gate $v$ and exponent
vector $\alpha$ of total degree $O_p(d)$, write $C(v,\alpha)$ for its
coefficient modulo $p$. Constants, variables, sums and differences are
immediate. At a product gate with predecessors $v_1,v_2$, use
\[
 C(v,\alpha)=\sum_{0\le\beta\le\alpha}
           C(v_1,\beta)C(v_2,\alpha-\beta).
\]
Evaluate the summands and the two factors sequentially. Each recursive
frame stores its gate index, a constant number of exponent vectors,
the current summation vector and field elements. Its space is
$O_p(d\log(d+2)+\log(L+2))$. Gate indices decrease along every recursive
path, so at most $O(L)$ frames are active. The resulting space bound is
\begin{equation}\label{fg:eq:coefficient-space}
 O_p\bigl(Ld\log(d+2)+L\log(L+2)\bigr).
\end{equation}
This procedure is exact even if intermediate gates have very large
formal degree: a term contributing to $x^\alpha$ at a product uses
only componentwise smaller nonnegative exponents. Arithmetic remains
modulo the fixed prime. Recomputed subcircuits cost time but no table
of coefficients. A circuit for $g=F^{p-1}$ has $O_p(L)$ gates, so this
also computes each $c_\beta$.

Next compute one coefficient of $\Delta$. For $|\gamma|=pD$, enumerate
the ordered $p$-tuples of degree-$D$ monomials whose exponents sum to
$\gamma$. Accumulate their products of canonical coefficients
$\widehat c_\beta$ modulo $p^2$. Subtract
$\widehat c_{\gamma/p}^{p}$ if every coordinate of $\gamma$ is divisible
by $p$, and subtract zero otherwise. The resulting canonical residue
in $\{0,\ldots,p^2-1\}$ is divisible by $p$; divide it by $p$ and reduce
modulo $p$. This is exactly \eqref{fg:eq:delta}. Since $p$ is fixed,
the tuple counters add only $O_p(d\log(d+2))$ bits to
\eqref{fg:eq:coefficient-space}. The computation uses the coefficients
of the collected polynomial $g$, so no uncollected parse-tree
multiplicity is mistaken for a monomial coefficient.

Let $M$ be the matrix of $T$ in the degree-$D$ monomial basis. Its row
$\alpha$, column $\beta$ entry is
\begin{equation}\label{fg:eq:matrix-entry}
 M_{\alpha\beta}
 =[x^{p\alpha+(p-1)\mathbf1-\beta}]\Delta,
\end{equation}
with value zero if a coordinate is negative. The target has total
degree $pD$, so the preceding procedure supplies this entry in the
same space. The matrix itself is never stored. Its initial vector has
component $c_\beta$, and its target functional selects row $\tau$.

For $e\ge0$, compute an entry $P(e,\alpha,\beta)=(M^e)_{\alpha\beta}$
by balanced recursion. At $e=0$ return the Kronecker delta, and at
$e=1$ use \eqref{fg:eq:matrix-entry}. For $e\ge2$, put
$a=\lfloor e/2\rfloor$ and $b=e-a$, and use
\[
 P(e,\alpha,\beta)=\sum_{|\eta|=D}
                   P(a,\alpha,\eta)P(b,\eta,\beta).
\]
Again evaluate both factors and all summands sequentially. For $e<B$
the recursion has depth $O_p(d)$, and each frame uses
$O_p(d\log(d+2))$ bits. Its total stack is
$O_p(d^2\log(d+2))$. At a leaf, the coefficient procedure adds its
own stack from \eqref{fg:eq:coefficient-space}; these space costs add.

For each $j=0,\ldots,B-1$, evaluate
\[
 a_j=\sum_{|\beta|=D}P(j,\tau,\beta)c_\beta.
\]
Output $j+1$ at the first nonzero value, and infinity if all $B$ tests
vanish. The outer counters use only $O_p(d\log(d+2))$ further bits.
Combining the bounds and $d=O(L)$ proves the asserted working space.
The finite answer has $O_p(d)$ bits. A binary-specified query can first
be compared with $B$ before this same algorithm is used.

Every loop has an explicit finite endpoint, and recursive gate indices
or power exponents decrease. On a nonhomogeneous off-promise input,
the degree-$D$ coefficients of $F^{p-1}$ define the vector used in the
same finite computation. Thus the algorithm halts without recognizing
the geometric promise. Linearity here is over $\Fp$; no semilinear
operator over an extension field has been silently substituted.
\end{proof}

The proof also gives a deterministic time bound
$2^{O_p(L^2\log(L+2))}$. The coefficient recursion has depth $O(L)$
and at most $2^{O_p(d\log(d+2))}$ children per frame, and balanced
matrix powering has $2^{O_p(d^2)}$ leaves. Repeating for $B$ indices
does not change the stated bound. The space conclusion is stronger than
this exponential-time conclusion. It is an upper bound: no matching
hardness classification for infinite height follows from it.

\begin{corollary}[Polynomial time in sparse or fixed-dimensional input]
\label{fg:cor:elementary-p}
Fix $p$. Every fixed finite exact-height predicate and every fixed
height-threshold predicate is decidable in deterministic polynomial time
for an explicitly listed sparse equation in the geometric promise above.
For a fixed number of variables, the full height is computable in
deterministic polynomial time from a division-free circuit.
\end{corollary}
\begin{proof}
For a fixed queried height, only a fixed finite list of coefficients
$b_{p^j}$ and congruences modulo $p^j$ is needed, by
\eqref{fg:eq:all-height} and \eqref{fg:eq:exact-height}. If the sparse
equation has $t$ terms, expand its $(p^j-1)$-st power by enumerating
ordered tuples of terms. There are $t^{p^j-1}$ tuples, a polynomial
number because $p,j$ are fixed. Add exponent vectors in binary, retain
the required diagonal exponent, and accumulate coefficients modulo
$p^j$. First collect equal exponent vectors modulo $p$, remove zero
coefficients, and lift the resulting homogeneous coefficients to
$\{0,\ldots,p-1\}$; this fixes an integral homogeneous lift even if
the input list originally had repeated or canceling terms.
A promised smooth input has every variable present, so the
number of variables is bounded by the length of the explicit input.
All these operations therefore have polynomial bit cost.

If $d$ is fixed, truncate the circuit at degree $d$ in its $d$ variables.
The truncated ring has the constant number $\binom{2d}{d}$ of basis
monomials. Evaluate its gates modulo $p$, then form the collected
$g=F^{p-1}$ and the constant-size operator $T$ in
\eqref{fg:eq:finite-operator}. All relevant degree bounds and the
matrix size $B$ are constant. Test its first $B$ target coefficients
by ordinary matrix-vector iteration. The total cost is polynomial in
the circuit length. Both algorithms can reject the elementary size
failures above and halt with arbitrary output outside the promise.
\end{proof}

\subsection{The cohomological and zeta-function comparison}
There is an independent finite method through the zeta function.
For $n=d-2>0$, rational Witt cohomology vanishes outside degrees zero
and $n$. To see this, use the truncation sequences
\[
 0\longrightarrow\OO_X\longrightarrow W_s\OO_X
   \longrightarrow W_{s-1}\OO_X\longrightarrow0
\]
and the coherent vanishings. Induction gives the intermediate
vanishings, while $H^0(X,W_s\OO_X)=W_s(\Fp)$.
The degree-zero and degree-$n$ transition maps are surjective, the
latter by $H^{n+1}(X,\OO_X)=0$. Their inverse systems satisfy the
Mittag--Leffler condition, so passage to Witt vectors introduces no
additional cohomology in adjacent degrees.

The rational Dieudonn\'e module of $\Phi_X$ is
$H^n(X,W\OO_X)[1/p]$
\citep[II, Corollary~4.3, p.~110]{fg-artin-mazur}.
It is zero at infinite height and has dimension $h$ at finite height;
this dichotomy is stated explicitly in
\citet[\S3, p.~6 of the preprint]{fg-yobuko}.
The comparison with the slope interval $[0,1)$ in crystalline cohomology
is \citet[II, Corollary~3.5, (3.5.3), p.~616]{fg-illusie}; see also
\citet[Theorem~1.1 and (6.2)]{fg-bbe}. Consequently
\begin{equation}\label{fg:eq:zeta-infinity}
 h(X)=\infty\iff\bigl((1-T)Z(X,T)\bigr)^{<1}=1.
\end{equation}
The factor $1-T$ removes degree-zero cohomology. All remaining
slope-less-than-one factors come from the top Witt group, with one
parity, so they cannot cancel each other.

For an exact algorithm, expand the input circuit modulo the ideal of
all monomials of total degree greater than $d$. At most
$\binom{2d}{d}\le4^d$ monomials occur; truncation is a ring
homomorphism and is valid at every gate regardless of its formal
degree. Partition projective space by the first nonzero coordinate,
set it equal to one, and specialize the earlier coordinates to zero.
This gives $d$ disjoint affine hypersurfaces with at most $d-1$
variables and degree at most $d$. Zeta functions multiply across the
partition. Zero and nonzero constant equations give affine spaces and
empty sets and are handled directly. For each other chart,
\citet[Theorem~37, p.~29 of the preprint]{ff-lauder-wan} give the
deterministic bit bound
\[
 \widetilde O\left(2^{13r^2}a^{3r+7}
                   e^{3r^2+9r}p^{2r+4}\right)
\]
for the full zeta function of a degree-$e$ polynomial in $r$ variables
over $\F_{p^a}$, without a smoothness hypothesis. With $a=1$, $r<d$,
$e\le d$, this is $2^{O_p(d^2\log(d+2))}$.

Reduce $(1-T)Z(X,T)=A(T)/B(T)$ to coprime integer polynomials with
constant coefficient one. Products, exact gcds and divisions have
polynomial bit cost in the explicit zeta output. Coprimeness over
$\Q$ persists over $\Q_p$. For
$C(T)=1+\sum_i c_iT^i$, all reciprocal roots have valuation at least
one exactly when $p^i\mid c_i$ for every $i$. One direction follows
from elementary symmetric functions; for the other, $C(U/p)$ has
integral coefficients, so its reciprocal polynomial is monic integral
and its roots are integral. Apply this test to both $A$ and $B$ in
\eqref{fg:eq:zeta-infinity}. This gives the same deterministic
exponential-time bound, with an independent cohomological explanation
of what the infinite-height predicate detects.

The finite-height estimate
$h\le\dim H^{n-1}(X,\Omega_X^1)+1$ of
\citet[Corollary~2.3, p.~100]{fg-vdgk} gives another geometric stopping
bound. The operator argument above supplies a bound directly from the
degree and variable count, while controlling the storage needed for a
succinct input.
